\documentclass[../main.tex]{subfiles}

\begin{document}

\chapter{Literatur}
\vspace{-0.5cm}
\noindent\textit{\small Verfasst von Vladana}
\vspace{0.5cm}

\section{Untersuchungsgebiet Ötztaler Alpen}

Die Region wird geologisch durch den Ötztal-Stubai-Kristallinkomplex geprägt, der überwiegend aus polymetamorphen Gesteinen wie Paragneisen, Glimmerschiefern und Orthogneisen besteht \autocite{dufresneRiverdammingLateQuaternaryRockslides2018, geistUpperOtzValley2022, kerschnerForelandCentralAlps}.
Die Wildspitze stellt mit 3768\,m die höchste Erhebung der Ötztaler Alpen dar und ist gleichzeitig der zweithöchste Gipfel Österreichs.
Insgesamt liegen in dem Gebiet viele Gipfel über 3000\,m \autocite{geistUpperOtzValley2022, aulitzkyMountainHazardGeomorphology1994, kerschnerForelandCentralAlps, bachnetzerBachnetzerTourismusentwicklungIm2013}.
Zudem ist die Landschaft durch pleistozäne Gletschererosion geprägt, die durch U-förmige Trogtäler, Kare oder Rundhöcker definiert wird.
Auch heute ist das Gebiet noch eine der am intensivsten vergletscherten Regionen in Österreich. Der Hintereisferner, Vernagtferner und der Gurgler Ferner gehören zu den bekanntesten und zu den am besten untersuchten Gletschern der Ostalpen \autocite{geistUpperOtzValley2022, aulitzkyMountainHazardGeomorphology1994, kerschnerForelandCentralAlps}.

\begin{figure}[H]
    \centering
    \includegraphics[width=0.7\textwidth]{img/literatureOne.png}
    \caption{Untersuchungsgebiet Ötztaler Alpen}
    \label{fig:literatureOne}
\end{figure}

Die bekannte Mumie von „Ötzi\" wurde ebenfalls in den Ötztaler Alpen gefunden. Dieser archäologische Fund ist etwa 5300 Jahre alt und belegt damit die frühe Anwesenheit und Einflussnahme des Menschen im Hochgebirge \autocite{geistUpperOtzValley2022, bachnetzerBachnetzerTourismusentwicklungIm2013}. Die aktive Landnutzung begann vor etwa 3500 Jahren. Sie umfasste umfangreiche Rodungen zur Ausweitung von Weideflächen. Das senkte die Waldgrenze dauerhaft ab. Mit den Rofenhöfen (2014\,m) und anderen Weilern befinden sich auch in hohen Lagen noch dauerhafte Siedlungsfläche in den hinteren Teilen der Täler \autocite{geistUpperOtzValley2022}.

Ursprünglich dominierten die Subsistenzwirtschaft und Viehhaltung, insbesondere die Schafzucht \autocite{stottenLockinsCommunityResilience2021, geistUpperOtzValley2022}. Diese wurde jedoch seit dem 19. Jahrhundert fast komplett durch den Tourismus ersetzt. Heute ist die Region ein internationales Zentrum für den Winter- und Sommersport, geprägt durch Bergsteigerdörfer als auch Massentourismus \autocite{stottenLockinsCommunityResilience2021, geistUpperOtzValley2022, aulitzkyMountainHazardGeomorphology1994, fischerGlaciersSnowSki2011, bachnetzerBachnetzerTourismusentwicklungIm2013}.

Als aktuelle Herausforderung gilt auch in den Ötztaler Alpen der Klimawandel mit einem Anstieg der Jahresmitteltemperatur um 1,2\,°C \autocite{stottenLockinsCommunityResilience2021, geistUpperOtzValley2022, fischerGlaciersSnowSki2011}. Gletscherrückgang und auftauender Permafrost erschwert die Instandhaltung der Infrastruktur, wie etwa Hütten, Wege oder Liftanlagen \autocite{stottenLockinsCommunityResilience2021, abadMassMovementSusceptibility2022, geistUpperOtzValley2022, fischerGlaciersSnowSki2011}. Auch die Nutzung von Kunstschnee, aufgrund des geringeren festen Niederschlags, stellt neue ökologische Probleme dar \autocite{stottenLockinsCommunityResilience2021, geistUpperOtzValley2022}. Zudem steigt das Gefahrenpotenzial, welches sich bereits in der Vergangenheit in Katastrophen wie Hochwasser, Murgängen und Gletscherseeausbrüchen äußerte \autocite{abadMassMovementSusceptibility2022, geistUpperOtzValley2022, aulitzkyMountainHazardGeomorphology1994}.

\section{Massenbewegungen in den Ötztaler Alpen}

Rutschungen bzw. Bergstürze stellen eine weitere Gefahr dar, mit der sich in diesem Projekt auseinandergesetzt wurde. Ein großes Ereignis in den Ötztaler Alpen war der Köfels-Bergsturz, der als die größte Massenbewegung in kristallinem Gestein im Alpenraum bekannt ist \autocite{hermannsExamplesMultipleRockslope2006, pragerGeologicalControlsSlope2009}. Mit einem Volumen von über 3,2\,km\textsuperscript{3} ereignete sich der Felssturz im frühen Holozän. Die Schuttmassen lösten sich von einem nach Osten gerichteten Hang am Wenderkogel. Sie füllten das Ötztal auf einer Länge von mehreren Kilometern auf, staute die Ötz und schuf das Becken von Längenfeld sowie das Tumpener Becken \autocite{pragerGeologicalControlsSlope2009}. Die Ötz schnitt sich später durch die Ablagerungen und so bildete sich die Maurach-Schlucht. Auch in anderen Tälern der Ötztaler Alpen kam es zu Bergstürzen, die prägend für die Geomorphologie waren \autocite{hermannsExamplesMultipleRockslope2006, pragerGeologicalControlsSlope2009}.

Massenbewegungen können in unterschiedliche Typen eingeteilt werden. Fels- oder Bergstürze wie eben beschrieben sind u.\,a. gesteinsabhängig und finden in kompetenten Gesteinen, wie Orthogneisen statt. Sie resultieren meist in schnellen Ereignissen. Über längere Zeiträume können Rutschungen und Kriechbewegungen stattfinden. Für sie sind inkompetente Gesteine wie Paragneise typisch \autocite{pragerGeologicalControlsSlope2009}. Ein Beispiel dafür ist die Wurzberg-Rutschung im Untersuchungsgebiet. Sie ist eine langsam kriechende Felsgleitung in inkompetenten Paragneisen südlich des Köfels-Abbruchgebiets. InSAR-Analysen zeigten dort Verschiebungsraten von 0,5 bis 1,5\,cm pro Jahr \autocite{pragerGeologicalControlsSlope2009}.

Verschiedene Faktoren und Ursachen begünstigen Rutschungen in den Ötztaler Alpen. Zum einen trägt die postglaziale Entwicklung dazu bei. Vorher vergletscherte Hänge werden nun freigelegt, sodass Risse entstehen, wachsen und Spannungsänderung hervorrufen können \autocite{pragerGeologicalControlsSlope2009}. Hänge werden instabiler und anfälliger für klimatische Trigger wie Extremniederschläge, Verwitterung und Permafrostdegradation \autocite{casagliLandslideDetectionMonitoring2023, hermleTimelyPredictionPotential2021, hermlePerformanceTestingOptical2022}.

Insbesondere der Wandel des Klimas trägt zu einer erhöhten Gefahr bei \autocite{hermleTimelyPredictionPotential2021, hermlePerformanceTestingOptical2022}, was betont, warum auch die SSP-Szenarien modellrelevant sind. Außerdem unterstreichen auch die historischen Ereignisse die Relevanz der Vorhersage von Rutschungen in den Ötztaler Alpen, besonders in Siedlungsnähe \autocite{casagliLandslideDetectionMonitoring2023, hermleTimelyPredictionPotential2021}.

\section{Fernerkundungsbasierte Detektion und Monitoring von Landslides}

Optische Verfahren gelten als eine der Hauptmethoden der Fernerkundung zur Untersuchung von Rutschungen \autocite{kaabRemoteSensingGlacier2005, scaioniRemoteSensingLandslide2014}. Multispektrale Satellitendaten (z.\,B. Sentinel-2) bieten eine hohe zeitliche Auflösung (Wiederholungsrate von bis zu 3 Tagen in mittleren Breiten), sind frei verfügbar und liefern Daten mit Auflösungen von 10\,m, 20\,m und 60\,m. Damit ermöglichen sie eine kontinuierliche Überwachung von großen Gebieten und die Erstellung multitemporaler Inventare. Die Nutzbarkeit ist stark von meteorologischen Bedingungen, wie der Wolkenbedeckung, und Tageslicht abhängig \autocite{vanwykdevriesDetectionSlowmovingLandslides2024, kaabRemoteSensingGlacier2005, scaioniRemoteSensingLandslide2014, strzabalaIdentificationLandslidePrecursors2024}.

Mithilfe dieser Daten können Indizes als Indikatoren für Bodeninstabilitäten genutzt werden, da Rutschungen oft die Pflanzendecke stören \autocite{piacentiniStatisticalAnalysisAssessing2012, vanwykdevriesDetectionSlowmovingLandslides2024}. Der NDVI (Normalized Difference Vegetation Index) wird häufig genutzt, da ein Rückgang des NDVI-Wertes als Proxy für den Verlust der Vegetationsbedeckung infolge einer Rutschung dient \autocite{vanwykdevriesDetectionSlowmovingLandslides2024, scaioniRemoteSensingLandslide2014, strzabalaIdentificationLandslidePrecursors2024}.

Mithilfe der Change Detection Methode kann außerdem die zeitliche Entwicklung von Hängen überwacht werden. Dabei werden Bilder verschiedener Zeitpunkte (bspw. vor und nach einem Ereignis) verglichen, um radiometrische oder strukturelle Unterschiede zu identifizieren \autocite{kaabRemoteSensingGlacier2005, scaioniRemoteSensingLandslide2014, strzabalaIdentificationLandslidePrecursors2024}.

Ein weiterer Ansatz bei dem optische Fernerkundungsdaten zum Einsatz kommen ist das Feature Tracking. Dabei wird die horizontale Verschiebung von Mustern zwischen Bildpaaren gemessen und kann für die Erkennung langsamer Rutschungen, die keine offensichtlichen Vegetationsverluste verursachen, angewendet werden \autocite{vanwykdevriesDetectionSlowmovingLandslides2024}. Objektbasierte Bildanalyse (OBIA) ist ein ähnliches Verfahren, doch anstatt einzelne Pixel zu betrachten, gruppiert OBIA benachbarte Pixel zu Objekten oder Segmenten. Diese Segmente werden dann nach ihrem farblichen Spektrum, ihrer Form, Textur, Größe und ihrem räumlichen Kontext klassifiziert \autocite{jinAccurateLandslideIdentification2022, scaioniRemoteSensingLandslide2014}.

Neben multispektralen Daten nutzt InSAR (z.\,B. Sentinel-1) die Phasendifferenz zwischen zwei oder mehr Radaraufnahmen, um Bodenverformungen mit Millimeterpräzision zu messen \autocite{kaabRemoteSensingGlacier2005, scaioniRemoteSensingLandslide2014, strzabalaIdentificationLandslidePrecursors2024}. Da Radarwellen Wolken durchdringen können, ermöglichen diese Satelliten eine kontinuierliche Überwachung unabhängig von Wetterbedingungen oder Tageszeit. Es gibt ebenfalls viele frei verfügbare Daten mit etwa sechs Tagen Wiederholrate \autocite{kaabRemoteSensingGlacier2005, scaioniRemoteSensingLandslide2014, strzabalaIdentificationLandslidePrecursors2024}.

DInSAR-Methoden (Differential) sind für das Monitoring langsam kriechender Erdrutsche geeignet, da sie oberflächennahe Deformationen mit hoher Präzision über lange Zeiträume erfassen können. Verfahren wie PS-InSAR (Persistent Scatterer) und SBAS-InSAR (Small Baseline Subset) verbessern die zeitliche Analyse, indem sie stabile Rückstreuzentren identifizieren bzw. interferometrische Bildpaare mit kleinen Basislinien kombinieren und dadurch atmosphärische Störungen sowie Dekorrelationseffekte reduzieren \autocite{kaabRemoteSensingGlacier2005, scaioniRemoteSensingLandslide2014, strzabalaIdentificationLandslidePrecursors2024}. So lassen sich Verschiebungsraten konsistent über mehrere Jahre quantifizieren. Zudem ermöglicht die Zeitreihenanalyse die Identifikation potenzieller Vorläufersignale, durch die Detektion von Beschleunigungen in den Deformationsraten Wochen bis Monate vor einem möglichen Kollaps. Ein signifikanter Rückgang der interferometrischen Kohärenz kann ebenfalls als Signal für strukturelle Veränderungen der Bodenoberfläche gewertet werden, da er auf Materialumlagerungen oder zunehmende Oberflächenrauigkeit hinweist, noch bevor größere Verschiebungen visuell sichtbar werden \autocite{scaioniRemoteSensingLandslide2014, strzabalaIdentificationLandslidePrecursors2024}.

Satellitengestützte Daten werden oft für großräumige Screenings genutzt, während GB-InSAR (ground-based) für das lokale Monitoring in Echtzeit eingesetzt wird \autocite{vanwykdevriesDetectionSlowmovingLandslides2024, scaioniRemoteSensingLandslide2014, strzabalaIdentificationLandslidePrecursors2024}. In Frühwarnsystemen können automatische Alarmmeldungen ausgelöst werden, wenn die gemessenen Geschwindigkeiten oder Geschwindigkeitszunahmen vordefinierte Schwellenwerte überschreiten \autocite{scaioniRemoteSensingLandslide2014, strzabalaIdentificationLandslidePrecursors2024}.

Auch in den Ötztaler Alpen wurden langsame Rutschungen in der Vergangenheit mit InSAR-Verfahren überwacht \autocite{pragerGeologicalControlsSlope2009, kaabRemoteSensingGlacier2005}. Diese Ansätze sind jedoch eher Ereignis- oder Monitoringorientiert. Sie stellen aktuelle Trends dar und dienen kurzfristigen Prognosen \autocite{kaabRemoteSensingGlacier2005, strzabalaIdentificationLandslidePrecursors2024}. Die vorliegende Projektarbeit hingegen zielt auf langfristige Vorhersagen. Multisource Datenfusion dient als geeignetere Basis für Zukunftsszenarien als Fernerkundungsdaten \autocite{woodLandslideInventoriesClimate2015, jinAccurateLandslideIdentification2022, kaabRemoteSensingGlacier2005, strzabalaIdentificationLandslidePrecursors2024}.

\section{Landslide Susceptibility Mapping und Machine Learning}

Im Landslide Susceptibility Mapping (LSM) hat sich die Multisource-Datenfusion zudem durchgesetzt. Es werden verschiedene Faktoren mit topographischen Daten aus digitalen Geländemodellen und teils auch Fernerkundungsdaten kombiniert, um die Genauigkeit der Vorhersagen zu maximieren und Fehlinterpretationen zu minimieren \autocite{jinAccurateLandslideIdentification2022, kaabRemoteSensingGlacier2005, yuLandslideSusceptibilityMapping2023}. So kann bspw. OBIA mit maschinellem Lernen kombiniert werden, um zwischen Rutschungsflächen und ähnlichen Merkmalen wie Straßen oder gerodeten Waldflächen zu unterscheiden \autocite{jiaNonlandslideSampleLandslide2025, wenSHAPenhancedXGBoostModel2024, jinAccurateLandslideIdentification2022, strzabalaIdentificationLandslidePrecursors2024}.

Statistische Verfahren zur Vorhersage von Rutschungsempfindlichkeit sind Logistische Regression oder Weight of Evidence, ein nicht-lineares statistisches Modell, das auf dem Bayes-Theorem basiert und die Beziehung zwischen Einflussfaktoren und Rutschungsereignissen über Wahrscheinlichkeiten definiert \autocite{piacentiniStatisticalAnalysisAssessing2012, goetzEvaluatingMachineLearning2015}. Hingegen können Baumbasierte ML-Modelle wie Random Forest oder Gradient-Boosting Decision Trees komplexe, nicht-lineare Beziehungen besser erfassen. Insbesondere Random Forest wird in vielen Studien als eines der leistungsfähigsten Modelle hervorgehoben, da es durch die Aggregation vieler Entscheidungsbäume stabil gegenüber Ausreißern ist und das Risiko von Overfitting minimiert \autocite{jiaNonlandslideSampleLandslide2025, goetzEvaluatingMachineLearning2015, huangModellingLandslideSusceptibility2024, yuLandslideSusceptibilityMapping2023}.

In der binären Klassifikation von Rutschungsereignissen wird zwischen Presence-Daten, also dokumentierte Rutschungen, und Absence-Daten unterschieden. Da Absence-Daten im geomorphologischen Kontext nicht eindeutig bestimmbar sind, werden Pseudo-Absences verwendet, also generierte Nicht-Rutschungsproben außerhalb bekannter Ereignisflächen. Die Qualität der Modellvorhersage hängt von der Auswahl dieser Referenzpunkte ab: Eine zufällige Generierung kann instabile Hangbereiche falsch als stabil klassifizieren und führt zu verzerrten Entscheidungsgrenzen. Semi-überwachte Ansätze wirken diesem Problem durch eine iterative Vorgehensweise entgegen. Dabei wird ein Vorabmodell trainiert, um anschließend Pseudo-Absences gezielt in Bereichen mit sehr niedrig vorhergesagter Suszeptibilität auszuwählen. Dadurch kann die Trennschärfe zwischen stabilen und instabilen Hangbereichen erhöht werden \autocite{piacentiniStatisticalAnalysisAssessing2012}.

Bei der Form der räumlichen Datenverteilung kann in raster- und vektorbasierte Ansätze unterschieden werden \autocite{scaioniRemoteSensingLandslide2014}. Rasterbasierte Modelle unterteilen das Gebiet in gleichgroße Pixel, denen jeweils Prädiktorwerte zugewiesen werden \autocite{wenSHAPenhancedXGBoostModel2024, huangModellingLandslideSusceptibility2024}. Vektorbasierte Ansätze wie Slope Units orientieren sich an morphologisch und hydrologisch abgegrenzten Hangsegmenten. Innerhalb dieser Einheiten werden die Prädiktorvariablen aggregiert, wodurch geomorphologische Prozesszusammenhänge konsistenter abgebildet werden \autocite{wenSHAPenhancedXGBoostModel2024, huangModellingLandslideSusceptibility2024}.

Deep-Learning-Modelle, wie Convolutional Neural Networks oder U-Net segmentieren bspw. Waldveränderungen pixelweise und können somit aktuelle Rutschungsereignisse aus Bilddaten automatisch erkennen. In Studien wurden sehr genaue Ergebnisse bei der Kartierung von Gefahrengebieten erzielt, doch die Modelle benötigen dafür eine große Menge an Trainingsdaten, die oft nicht flächendeckend verfügbar sind \autocite{jinAccurateLandslideIdentification2022, pyoGeneralizationUNetSemantic2022}. Aufgrund ihres Fokus auf die Detektion visueller Merkmale in aktuellen Bildern eignen sie sich hauptsächlich für das Monitoring und sind begrenzt in der Anwendung für langfristige, szenariobasierte Projektionen, die auf statistischen Korrelationen beruhen \autocite{mondiniShortLongTerm2025, jinAccurateLandslideIdentification2022, pyoGeneralizationUNetSemantic2022}.

\end{document}